\section{Method}
\label{sec:method}

We propose \emph{Typed Composition Search (TCS)}, a formulation of tool routing that separates semantic reasoning from compositional reasoning.
Instead of predicting executable tools directly, a language model predicts the entity types involved in a query, while deterministic graph search recovers a valid tool composition.
LLMs perform semantic reasoning over entity types; graph search performs compositional reasoning over valid transformations.

\subsection{Problem Formulation}
\label{sec:problem-def}

Existing approaches formulate tool routing as direct tool selection from a catalog.
Given a natural language query $q$, a model selects one or more tools:

\begin{equation}
    g(q) \rightarrow \{t_1, \ldots, t_k\} \subseteq \mathcal{T}
\end{equation}

This formulation requires the model to simultaneously identify relevant tools and infer a valid execution order, without any guarantee that the resulting composition is structurally valid.

We instead formulate routing over typed entities.

Let $\mathcal{E} = \{e_1, e_2, \ldots, e_n\}$ denote the set of entity types in a domain (e.g., \texttt{Pod}, \texttt{Namespace}, \texttt{Deployment} in Kubernetes).
Let $\mathcal{T} = \{t_1, t_2, \ldots, t_m\}$ denote the available tools, where each tool is a typed transformation

\begin{equation}
    t_k : e_i \rightarrow e_j,
    \qquad
    e_i,e_j \in \mathcal{E}.
    \label{eq:tool-transform}
\end{equation}

That is, tool $t_k$ consumes an entity of type $e_i$ and produces an entity of type $e_j$.

Given a query, the objective is to recover an ordered tool sequence
$[t_1,\ldots,t_l]$
that transforms a source entity into the desired target entity.
The sequence must be compositionally valid: the output type of every tool must match the input type of the next.

\subsection{Typed Composition Graph}
\label{sec:typed-graph}

The tool registry is represented as a directed graph

\[
G=(\mathcal{E},\mathcal{T}),
\]

where entity types are nodes and tools are directed edges.
An edge exists from $e_i$ to $e_j$ whenever a tool
$t_k:e_i\rightarrow e_j$
is available.

Multiple tools may connect the same pair of entity types, and tools with multiple input or output types generate multiple edges.

In this representation,

\begin{itemize}
\item a valid tool invocation corresponds to a graph edge;
\item a valid multi-step workflow corresponds to a graph path.
\end{itemize}

For example, consider a Kubernetes tool registry:

\begin{itemize}
    \item \texttt{list\_pods}: \texttt{Namespace} $\rightarrow$ \texttt{Pod}
    \item \texttt{get\_pod\_logs}: \texttt{Pod} $\rightarrow$ \texttt{PodLogs}
    \item \texttt{get\_pod\_node}: \texttt{Pod} $\rightarrow$ \texttt{Node}
\end{itemize}

A query such as

\begin{quote}
``Get the logs for pods in the production namespace.''
\end{quote}

implicitly specifies a path from
\texttt{Namespace}
to
\texttt{PodLogs}.
The corresponding tool sequence emerges directly from graph connectivity rather than being predicted explicitly by the language model.

The graph is constructed once from the typed tool registry and is independent of the language model.
No training, fine-tuning, embeddings, or learned routing policies are required.

\subsection{Typed Composition Search}
\label{sec:tcs}

Routing is decomposed into two independent stages.

First, the language model predicts the source and target entity types:

\begin{equation}
    f(q)
    \rightarrow
    (\hat e_{\mathrm{src}},
     \hat e_{\mathrm{tgt}})
    \label{eq:type-prediction}
\end{equation}

Second, deterministic graph search recovers the shortest valid composition between the predicted types.

The complete routing function is

\begin{equation}
    \pi(q)
    =
    \textsc{Path}
    \left(
        G,
        f_{\mathrm{src}}(q),
        f_{\mathrm{tgt}}(q)
    \right),
    \label{eq:routing-function}
\end{equation}

where $\textsc{Path}$ returns the shortest valid tool sequence connecting the predicted entity types.
Any shortest-path algorithm may be used; in our implementation we use breadth-first search because all graph edges are unweighted.

Since every returned sequence corresponds to a path in the typed graph, every tool composition is guaranteed to satisfy input-output compatibility.
Structurally invalid tool compositions cannot be generated.

This decomposition separates routing correctness into two independent components:

\begin{enumerate}
\item semantic correctness (whether the language model predicts the correct entity types);
\item compositional correctness (whether the graph contains a valid path between them).
\end{enumerate}

If no valid path exists, the system returns an empty result rather than generating an invalid tool composition.

\subsection{Graph-Constrained Reasoning}
\label{sec:graph-constraint}

The primary advantage of this formulation is not simply that there are fewer entity types than tools.
Rather, the graph itself constrains the reasoning space.

Once a target entity type has been predicted, reverse reachability eliminates every source type that cannot reach that target through any valid composition.

We quantify this structural constraint using the \emph{entity pruning} metric.

Let
$R^{-1}(e_{\mathrm{tgt}})\subseteq\mathcal{E}$
denote the set of entity types from which
$e_{\mathrm{tgt}}$
is reachable in the graph.

Entity pruning is defined as

\begin{equation}
    \mathrm{pruning}(e_{\mathrm{tgt}})
    =
    1-
    \frac{
        |R^{-1}(e_{\mathrm{tgt}})|
    }{
        |\mathcal{E}|
    }.
    \label{eq:entity-pruning}
\end{equation}

For example, a pruning value of 0.87 indicates that 87\% of possible source entity types are eliminated because they cannot reach the predicted target through any valid composition.

Importantly, this reduction depends only on graph topology.
It is independent of the language model and therefore applies equally to any type prediction model.

Because routing decomposes into type prediction followed by deterministic graph search, end-to-end recall also admits a simple decomposition:

\begin{equation}
    R
    =
    P_cR_c
    +
    (1-P_c)R_w,
    \label{eq:recall-decomposition}
\end{equation}

where $P_c$ denotes the probability of correct entity type prediction,
$R_c$ the routing recall when the predicted types are correct,
and $R_w$ the routing recall when they are incorrect.

This decomposition separates language model performance from graph quality and enables routing failures to be analyzed independently.